\documentclass{article} % For LaTeX2e
\usepackage{iclr2026_conference,times}

% Optional math commands from https://github.com/goodfeli/dlbook_notation.
\input{math_commands.tex}

\usepackage{hyperref}
\usepackage{url}


\title{Submodular Maximization with Optimal Hints}

% Authors must not appear in the submitted version. They should be hidden
% as long as the \iclrfinalcopy macro remains commented out below.
% Non-anonymous submissions will be rejected without review.

\author{Antiquus S.~Hippocampus, Natalia Cerebro \& Amelie P. Amygdale \thanks{ Use footnote for providing further information
about author (webpage, alternative address)---\emph{not} for acknowledging
funding agencies.  Funding acknowledgements go at the end of the paper.} \\
Department of Computer Science\\
Cranberry-Lemon University\\
Pittsburgh, PA 15213, USA \\
\texttt{\{hippo,brain,jen\}@cs.cranberry-lemon.edu} \\
\And
Ji Q. Ren \& Yevgeny LeNet \\
Department of Computational Neuroscience \\
University of the Witwatersrand \\
Joburg, South Africa \\
\texttt{\{robot,net\}@wits.ac.za} \\
\AND
Coauthor \\
Affiliation \\
Address \\
\texttt{email}
}

% The \author macro works with any number of authors. There are two commands
% used to separate the names and addresses of multiple authors: \And and \AND.
%
% Using \And between authors leaves it to \LaTeX{} to determine where to break
% the lines. Using \AND forces a linebreak at that point. So, if \LaTeX{}
% puts 3 of 4 authors names on the first line, and the last on the second
% line, try using \AND instead of \And before the third author name.

\newcommand{\fix}{\marginpar{FIX}}
\newcommand{\new}{\marginpar{NEW}}

%\iclrfinalcopy % Uncomment for camera-ready version, but NOT for submission.
\begin{document}


\maketitle

\begin{abstract}
Submodular maximization under cardinality constraints is a fundamental problem with broad applications in several areas such as machine learning, including data summarization, feature selection, and influence maximization. The classical framework assumes no prior knowledge about the possible solution sets, and the greedy algorithm provides the best possible approximation $(1-1/e)$ 
in this setting. We consider the submodular optimization problem with the prior knowledge of a partial optimal solution; we refer to this knowledge as optimal hints.  
There are two possible scenarios with cardinality constraint $\leq k$: 
one where some $\ell$ elements of optimal solution is given and another where  the "best" $\ell$ elements of some optimal solution is given. These scenarios are motivated by the idea that prior knowledge may become available due to domain knowledge from experts,  due to some more expensive prior computations, or as predictions of a learning algorithm. We show that the greedy strategy admits to 
$(1-1/e)f(OPT) + (1/e) f(OPT_{\ell})$-approximation for the first case and to $1 - ((k-\ell)/ek) (1 - 1/(k-\ell))^{-\ell}f(OPT)$-approximation for the second case. In addition to theoretical guarantees, we present preliminary experiments characterizing the situations where application of the greedy strategy, in practice, can lead to a close to optimal solution  in the presence of prior knowledge of partial optimal solutions. 

%We propose a technique in which the algorithm is given partial knowledge of the optimal solution—specifically, we are presenting with two different scenarios when running the greedy algorithm with a partial knowledge given. We consider the case where the hints are any $k-m$ elements and the best $k-m$ elements from an optimal solution $(OPT)$, for which we obtain a theoretical $(1-1/e)f(OPT) + (1/e) f(OPT_{m-k})$-approximation and $1 - (m/ek) (1 - 1/m)^{m-k}f(OPT)$-approximation respectively. These scenarios are motivated by the idea of using the optimal hints as an oracle-like model that could be used for a learning augmented algorithm.

%We present theoretical results showing that this technique enables stronger performance guarantees than those achievable in the classical framework. We also discuss the limits of this advantage and identify regimes where partial knowledge provides the greatest benefit.

\end{abstract}

\end{document}